\chapter{Preconditioning}
\label{ch:precond}

\chapterorient{The last two chapters built solvers and then bounded their speed by a
single number: the condition number of the operator. Conjugate gradients
(\cref{ch:cg}) crawls on a stretched bowl; GMRES (\cref{ch:gmres}) needs more steps
the more the spectrum spreads. This chapter is the cure. We do not build a faster
solver---we change the \emph{problem} the solver sees, multiplying it by an
approximate inverse $\mat{M}^{-1}\approx\mat{A}^{-1}$ that pulls the spectrum into a
tight cluster near $1$, so the very same iteration finishes in a fraction of the
steps. The spine of the chapter is the abstraction of \cref{ch:operators}: a
preconditioner is never a matrix you form, it is an \emph{operator you apply}---a
black-box function the solver calls and never looks inside. That one idea makes
preconditioning modular: the same GMRES drives a diagonal scaling, an incomplete
factorization, or a multigrid V-cycle without changing a line. We make the spectral
claim precise, sort out left/right/split placement, build the preconditioner once and
apply it each step (the constructed operator of \cref{ch:operators}), tour the
families, and end at the real goal: convergence that does not slow down as the mesh is
refined.}

\section{The motivation, made precise}
\label{sec:motivation}

We ended \cref{ch:cg} with a stark plot (\cref{fig:convplot}): the conjugate-gradient
error after $k$ steps is bounded by
\[
  \norm{\vx^\star-\vx_k}_{\mat{A}}
  \;\le\; 2\!\left(\frac{\sqrt{\kappa}-1}{\sqrt{\kappa}+1}\right)^{\!k}
  \norm{\vx^\star-\vx_0}_{\mat{A}},
  \qquad \kappa=\frac{\lambda_{\max}}{\lambda_{\min}},
\]
and GMRES carries an analogous spectral dependence. The number of iterations to reach
a fixed tolerance grows like $\sqrt{\kappa}$ for CG. The operators Palace assembles are
routinely ill-conditioned: refining the mesh, or putting a high-permittivity dielectric
next to a low one, drives $\kappa$ into the thousands and beyond. Raw CG and raw GMRES
are simply too slow on these systems. The fix is not a cleverer iteration---it is to
\emph{change the spectrum}.

\subsection{The idea in one line}

Suppose we have a second operator $\mat{M}$, close to $\mat{A}$ but \emph{cheap to
invert}, and consider solving the algebraically equivalent system
\begin{equation}
\label{eq:precond-system}
  \mat{M}^{-1}\mat{A}\,\vx \;=\; \mat{M}^{-1}\vb .
\end{equation}
It has the \emph{same solution} $\vx^\star$ as $\mat{A}\vx=\vb$---multiplying both sides
by the invertible $\mat{M}^{-1}$ changes nothing about the answer. But it can have a
completely different \emph{operator}: where $\mat{A}$ might have a spectrum spread across
$[\lambda_{\min},\lambda_{\max}]$ with $\kappa$ in the thousands, the new operator
$\mat{M}^{-1}\mat{A}$ has a spectrum we can hope to squeeze near $1$. In the ideal limit
$\mat{M}=\mat{A}$, the new operator is $\mat{M}^{-1}\mat{A}=\mat{I}$, whose condition
number is exactly $1$: a Krylov method solves $\mat{I}\vx=\mat{M}^{-1}\vb$ in a single
step. We will never reach that ideal cheaply (inverting $\mat{A}$ exactly is the problem
we are trying to avoid), but the closer $\mat{M}^{-1}$ approximates $\mat{A}^{-1}$, the
closer $\mat{M}^{-1}\mat{A}$ is to the identity, and the faster the iteration runs.

\begin{definition}[Preconditioner]
\label{def:precond}
A \emph{preconditioner} for the system $\mat{A}\vx=\vb$ is an operator $\mat{M}$,
nonsingular and \emph{cheap to apply the inverse of}, chosen so that
$\mat{M}^{-1}\mat{A}$ is better conditioned (its spectrum more clustered) than
$\mat{A}$. The action $\vr\mapsto\mat{M}^{-1}\vr$ is the \emph{preconditioner apply}.
\end{definition}

Two demands pull against each other in that definition, and a good preconditioner is a
truce between them. It must \emph{cluster the spectrum}---otherwise it does not reduce
the iteration count---and it must be \emph{cheap to apply}---otherwise the saved
iterations cost more than they were worth. The extremes are both useless:
$\mat{M}=\mat{I}$ is free to apply but clusters nothing (it does not precondition at
all); $\mat{M}=\mat{A}$ clusters perfectly but is as expensive to invert as the original
problem. Every practical preconditioner lives in between.

\subsection{The fundamental trade}
\label{sec:trade}

Make the bargain quantitative. Without preconditioning, a solve costs
\[
  (\text{iterations}) \times (\text{cost per iteration}),
\]
and the cost per iteration is dominated by one operator apply $\mat{A}\vp$. Preconditioning
changes \emph{both} factors, in opposite directions:

\begin{itemize}[nosep]
  \item \textbf{Iterations go down}---often dramatically---because the spectrum of
  $\mat{M}^{-1}\mat{A}$ is clustered and the Krylov bound improves.
  \item \textbf{Cost per iteration goes up}, because each step now performs one
  \emph{extra} apply: the preconditioner apply $\vz=\mat{M}^{-1}\vr$ on top of the
  operator apply $\mat{A}\vp$.
\end{itemize}

\begin{keyidea}
Preconditioning wins exactly when the product wins: the iteration count must fall by
more than the per-step cost rises. This happens when $\mat{M}^{-1}$ is
\emph{simultaneously} (i) a good approximation of $\mat{A}^{-1}$---so it genuinely
clusters the spectrum and slashes the iteration count---and (ii) cheap to apply---so the
extra per-step apply is a small surcharge, not a second full solve. A preconditioner
that clusters but is expensive, or one that is cheap but clusters nothing, loses the
trade. The art is finding an $\mat{M}^{-1}$ that is both.
\end{keyidea}

\Cref{fig:cost-trade} draws the bargain. As the preconditioner grows stronger (left to
right), the iteration count plunges, but the work spent \emph{inside} each preconditioner
apply climbs. Total time is the product, and it has a minimum somewhere in the middle: a
``sweet spot'' preconditioner strong enough to cut iterations sharply yet cheap enough
that each apply stays light. Too weak and the iterations dominate; too strong and the
per-step preconditioner apply dominates.

\begin{figure}[t]
\centering
\begin{tikzpicture}
  \begin{axis}[
      width=12.4cm, height=6.6cm,
      xlabel={preconditioner strength (cheap/weak $\to$ strong/expensive)},
      ylabel={cost (arbitrary units)},
      xmin=0, xmax=10, ymin=0, ymax=11,
      xtick=\empty,
      tick label style={font=\scriptsize}, label style={font=\footnotesize},
      legend pos=north west, legend cell align=left,
      legend style={font=\footnotesize, draw=simGray!40},
      grid=both, grid style={simGray!15},
      every axis plot/.append style={very thick}]
    % iterations: falling
    \addplot[simBlue, domain=0.3:10, samples=80]
      {1.0 + 9.0*exp(-0.55*x)};
    \addlegendentry{iterations (fall)}
    % per-iteration cost: rising
    \addplot[simRust, domain=0.3:10, samples=80]
      {0.6 + 0.34*x};
    \addlegendentry{cost per iteration (rise)}
    % total time = product, normalized down to fit
    \addplot[simGreen, domain=0.3:10, samples=120]
      {0.42*(1.0 + 9.0*exp(-0.55*x))*(0.6 + 0.34*x)};
    \addlegendentry{total time (product)}
    % mark the minimum of the green curve, near x ~ 3.4
    \node[circle, fill=simGreen, inner sep=1.6pt] at (axis cs:3.4,2.55) {};
    \node[font=\scriptsize, simGreen, anchor=south west] at (axis cs:3.55,2.7)
      {sweet spot};
    \draw[simGreen, densely dashed] (axis cs:3.4,0)--(axis cs:3.4,2.55);
  \end{axis}
\end{tikzpicture}
\caption[The fundamental preconditioning trade]{The trade at the heart of
preconditioning. A stronger preconditioner (rightward) clusters the spectrum more, so
the iteration count (blue) falls steeply---but each step now spends more work inside the
preconditioner apply, so the cost per iteration (rust) rises. Total solve time (green) is
the \emph{product} of the two, and it has a minimum in the middle: the sweet spot is a
preconditioner strong enough to cut iterations sharply yet cheap enough that each apply
stays light. At the left edge the iterations dominate (too weak); at the right edge the
per-step apply dominates (too strong, approaching a full inner solve). A good
preconditioner sits near the minimum.}
\label{fig:cost-trade}
\end{figure}

\section{What clustering buys: the spectrum picture}
\label{sec:spectrum}

The phrase ``better conditioned'' in \cref{def:precond} deserves a picture, because it
is the whole mechanism. \Cref{ch:cg} taught (\cref{sec:cluster}) that a Krylov method's
true convergence is governed not just by the extreme eigenvalues but by how the
\emph{whole} spectrum is laid out: a polynomial can be made small on a few tight clusters
with far fewer degrees than it needs to be small across a wide interval. A preconditioner
is, at bottom, a spectrum-reshaping device.

\Cref{fig:spectrum} shows the effect on a model problem. The eigenvalues of $\mat{A}$
(top) sprawl across two decades---this is the long, narrow bowl of \cref{ch:cg}, the
$\kappa$ in the thousands. The eigenvalues of $\mat{M}^{-1}\mat{A}$ (bottom), for a
decent preconditioner, huddle in a tight band around $1$. Same solution, same iteration,
utterly different convergence: the bottom spectrum is one a degree-handful polynomial
annihilates, so the Krylov method converges in a handful of steps.

\begin{figure}[t]
\centering
\begin{tikzpicture}
  \begin{axis}[
      width=13.2cm, height=4.8cm,
      xmode=log, log basis x=10,
      xmin=0.05, xmax=300,
      ymin=-0.6, ymax=2.2,
      ytick={0,1.5}, yticklabels={$\mat{M}^{-1}\mat{A}$, $\mat{A}$},
      y tick label style={font=\footnotesize},
      xlabel={eigenvalue (log scale)},
      label style={font=\footnotesize}, tick label style={font=\scriptsize},
      axis y line=left, axis x line=bottom,
      clip=false]
    % --- spectrum of A: spread across two decades ---
    \addplot[only marks, mark=*, mark size=2.0pt, simBlue]
      coordinates {(0.08,1.5)(0.15,1.5)(0.3,1.5)(0.7,1.5)(1.6,1.5)
                   (4,1.5)(9,1.5)(22,1.5)(55,1.5)(130,1.5)(240,1.5)};
    \node[font=\scriptsize, simBlue, anchor=west] at (axis cs:0.055,2.0)
      {spread: $\kappa\approx 3000$};
    % --- spectrum of M^{-1}A: clustered near 1 ---
    \addplot[only marks, mark=*, mark size=2.0pt, simGreen]
      coordinates {(0.62,0)(0.74,0)(0.83,0)(0.9,0)(0.97,0)(1.0,0)
                   (1.04,0)(1.1,0)(1.18,0)(1.3,0)(1.45,0)};
    \node[font=\scriptsize, simGreen, anchor=west] at (axis cs:1.6,0)
      {clustered: $\kappa\approx 2.3$};
    % a marker line at lambda = 1
    \draw[simGray, densely dashed] (axis cs:1,-0.55)--(axis cs:1,1.0)
      node[pos=0.0, below, font=\scriptsize, simGray] {$\lambda=1$};
  \end{axis}
\end{tikzpicture}
\caption[Preconditioning clusters the spectrum]{The spectral effect of preconditioning,
on a log scale. \emph{Top (blue):} the eigenvalues of the bare operator $\mat{A}$ sprawl
across two decades---a condition number in the thousands, the long narrow bowl that makes
CG crawl (\cref{ch:cg}). \emph{Bottom (green):} the eigenvalues of the preconditioned
operator $\mat{M}^{-1}\mat{A}$ collapse into a tight cluster around $\lambda=1$, with a
condition number near $2$. The system has the \emph{same solution}---only the operator
the Krylov method iterates against has changed. A polynomial that must be tiny across the
blue spread needs a high degree (many iterations); one that need only be tiny on the
green cluster needs a low degree (few iterations). Clustering near $1$ is exactly what a
good preconditioner buys.}
\label{fig:spectrum}
\end{figure}

\begin{notationbox}
``Clustered near $1$'' is the goal, and the $1$ is not arbitrary. If $\mat{M}^{-1}\mat{A}$
were clustered near any constant $c\neq0$, a rescaling would move the cluster to $1$, so
what matters is the \emph{spread} of the cluster, not its location. We say
$\mat{M}^{-1}\mat{A}\approx\mat{I}$ to mean its spectrum is a tight cluster; equivalently
$\mat{M}^{-1}\approx\mat{A}^{-1}$. The closer the approximation, the tighter the cluster,
the fewer the iterations---and the limit $\mat{M}^{-1}=\mat{A}^{-1}$ gives the single-step
solve.
\end{notationbox}

\section{Left, right, and split preconditioning}
\label{sec:lrs}

Equation~\eqref{eq:precond-system} put $\mat{M}^{-1}$ on the \emph{left} of $\mat{A}$.
That was one choice of three. The preconditioner can be applied on the left, on the
right, or split across both sides, and the three placements---algebraically equivalent in
their solution, but different in what the iteration actually minimizes---are a variant
axis (\cref{ch:operators}) every Krylov solver exposes.

\begin{description}[style=nextline,leftmargin=1.5em]
  \item[Left preconditioning.] Multiply the system on the left:
  \[
    \mat{M}^{-1}\mat{A}\,\vx \;=\; \mat{M}^{-1}\vb .
  \]
  The iteration runs on the operator $\mat{M}^{-1}\mat{A}$, and the residual it sees and
  drives to zero is the \emph{preconditioned residual}
  $\mat{M}^{-1}(\vb-\mat{A}\vx)=\mat{M}^{-1}\vr$, not the true residual $\vr$.
  \item[Right preconditioning.] Insert $\mat{M}^{-1}\mat{M}=\mat{I}$ between $\mat{A}$ and
  $\vx$, and solve for a transformed unknown $\vu=\mat{M}\vx$:
  \[
    \mat{A}\mat{M}^{-1}\,\vu \;=\; \vb, \qquad \vx=\mat{M}^{-1}\vu .
  \]
  The iteration runs on $\mat{A}\mat{M}^{-1}$, but---crucially---the residual it sees is
  $\vb-\mat{A}\mat{M}^{-1}\vu=\vb-\mat{A}\vx=\vr$, the \emph{true} residual. Right
  preconditioning leaves the quantity being minimized untouched.
  \item[Split preconditioning.] Factor the preconditioner $\mat{M}=\mat{M}_1\mat{M}_2$
  and place one factor on each side:
  \[
    \mat{M}_1^{-1}\mat{A}\mat{M}_2^{-1}\,\vu \;=\; \mat{M}_1^{-1}\vb,
    \qquad \vx=\mat{M}_2^{-1}\vu .
  \]
  When $\mat{A}$ is SPD and $\mat{M}=\mat{L}\mat{L}^{\top}$ is a (incomplete) Cholesky
  factorization, the split choice $\mat{M}_1=\mat{L}$, $\mat{M}_2=\mat{L}^{\top}$ keeps the
  preconditioned operator $\mat{L}^{-1}\mat{A}\mat{L}^{-\top}$ \emph{symmetric positive
  definite}---which is exactly what preconditioned CG needs (\cref{sec:pcg-tie}).
\end{description}

\begin{figure}[t]
\centering
\begin{tikzpicture}[font=\footnotesize, node distance=5mm]
  % --- LEFT ---
  \begin{scope}
    \node[font=\bfseries\small, simBlue] at (0,1.5) {left};
    \node[extern, minimum width=10mm] (lM) at (0,0.5) {$\mat{M}^{-1}$};
    \node[stage, minimum width=9mm, below=4mm of lM] (lA) {$\mat{A}$};
    \node[data, below=4mm of lA] (lx) {$\vx$};
    \draw[flow] (lx)--(lA); \draw[flow] (lA)--(lM);
    \node[font=\scriptsize, simGray, anchor=north, align=center] at (0,-1.55)
      {$\mat{M}^{-1}\mat{A}\vx=\mat{M}^{-1}\vb$\\[1pt]sees $\mat{M}^{-1}\vr$};
  \end{scope}
  % --- RIGHT ---
  \begin{scope}[xshift=4.7cm]
    \node[font=\bfseries\small, simGreen] at (0,1.5) {right};
    \node[stage, minimum width=9mm] (rA) at (0,0.5) {$\mat{A}$};
    \node[extern, minimum width=10mm, below=4mm of rA] (rM) {$\mat{M}^{-1}$};
    \node[data, below=4mm of rM] (ru) {$\vu$};
    \draw[flow] (ru)--(rM); \draw[flow] (rM)--(rA);
    \node[font=\scriptsize, simGray, anchor=north, align=center] at (0,-1.55)
      {$\mat{A}\mat{M}^{-1}\vu=\vb$\\[1pt]sees \emph{true} $\vr$};
  \end{scope}
  % --- SPLIT ---
  \begin{scope}[xshift=9.4cm]
    \node[font=\bfseries\small, simRust] at (0,1.5) {split};
    \node[extern, minimum width=10mm] (sM1) at (0,0.9) {$\mat{M}_1^{-1}$};
    \node[stage, minimum width=9mm, below=3mm of sM1] (sA) {$\mat{A}$};
    \node[extern, minimum width=10mm, below=3mm of sA] (sM2) {$\mat{M}_2^{-1}$};
    \node[data, below=3mm of sM2] (su) {$\vu$};
    \draw[flow] (su)--(sM2); \draw[flow] (sM2)--(sA); \draw[flow] (sA)--(sM1);
    \node[font=\scriptsize, simGray, anchor=north, align=center] at (0,-1.95)
      {$\mat{M}_1^{-1}\mat{A}\mat{M}_2^{-1}\vu=\mat{M}_1^{-1}\vb$\\[1pt]preserves symmetry};
  \end{scope}
\end{tikzpicture}
\caption[Three placements of the preconditioner]{The three placements of $\mat{M}^{-1}$,
read bottom-to-top as the operand flows through. \emph{Left:} $\mat{M}^{-1}$ multiplies
the whole system from the left; the iteration runs on $\mat{M}^{-1}\mat{A}$ and drives the
\emph{preconditioned} residual $\mat{M}^{-1}\vr$ to zero. \emph{Right:} $\mat{M}^{-1}$ sits
between $\mat{A}$ and the transformed unknown $\vu=\mat{M}\vx$; the iteration runs on
$\mat{A}\mat{M}^{-1}$ but sees the \emph{true} residual $\vr$---the reason right
preconditioning is the default for GMRES. \emph{Split:} one factor of
$\mat{M}=\mat{M}_1\mat{M}_2$ on each side preserves SPD-ness of the preconditioned
operator, which preconditioned CG requires. All three have the same solution $\vx^\star$;
they differ in what residual the convergence test reads.}
\label{fig:lrs}
\end{figure}

\subsection{Why the choice matters: the residual GMRES minimizes}

The three placements yield the same solution, so why prefer one? The answer is the
\emph{stopping test}. A Krylov method measures convergence by a residual norm, and the
three placements present it with different residuals.

\begin{keyidea}
\textbf{Right preconditioning preserves the true residual.} Left preconditioning makes
the iteration minimize $\norm{\mat{M}^{-1}\vr}$---the residual seen through the
preconditioner---so the stopping test is distorted by $\mat{M}$: a small
$\norm{\mat{M}^{-1}\vr}$ need not mean a small true $\norm{\vr}$ if $\mat{M}$ is
ill-scaled. Right preconditioning leaves the residual as the true $\vr=\vb-\mat{A}\vx$, so
GMRES's free residual readout (\cref{ch:gmres}) measures the quantity the engineer
actually cares about. This is why right preconditioning is the conventional default for
preconditioned GMRES; left is common for CG, where the preconditioned inner product
restores symmetry anyway (\cref{sec:pcg-tie}).
\end{keyidea}

There is one more reason the distinction is load-bearing, and it ties back to
\cref{ch:gmres}: for \emph{right}-preconditioned GMRES the iterate is recovered as
$\vx=\vx_0+\mat{M}^{-1}\mat{V}_m\vy$, applying $\mat{M}^{-1}$ once to the whole
correction---which assumes a \emph{single fixed} $\mat{M}$. The moment the preconditioner
varies per step, that reconstruction breaks, and FGMRES is the repair. We return to this
in \cref{sec:flexible}; for now, file the placement as a variant axis the construction
absorbs.

\section{The preconditioner is an operator you apply}
\label{sec:operator}

We have written $\mat{M}^{-1}$ throughout as if it were a matrix. It is not---and the gap
between ``$\mat{M}^{-1}$ the matrix'' and ``$\mat{M}^{-1}$ the operation'' is the single
most important idea in the chapter. \emph{The inverse is never formed.} What the solver
holds is a function: hand it a residual $\vr$, it returns a vector $\vz\approx\mat{M}^{-1}\vr$.
How it computes that vector---a diagonal divide, a triangular sweep, a full multigrid
cycle---is its own business, hidden behind the apply.

This is precisely the \emph{solver-as-operator} move of \cref{ch:operators}. There we saw
that a solver computes the action $\vv\mapsto\mat{M}^{-1}\vv$, which has the same shape as
the action $\vv\mapsto\mat{M}\vv$ of the operator it inverts---vector in, vector out---so
a solver can wear the very same \code{apply} interface as a plain operator. A
preconditioner \emph{is} a solver in this sense, and that is why a Krylov method needs
nothing from it but the apply.

\begin{keyidea}
\textbf{A preconditioner is an approximate inverse you apply; the solver never looks
inside.} The Krylov iteration is written against one interface---\code{apply(pc, r)}---and
asks exactly one thing of the preconditioner: turn this residual into an approximate
error, $\vz\approx\mat{M}^{-1}\vr$. It does not know, and cannot ask, \emph{how}. The same
GMRES drives a Jacobi diagonal-scale, an incomplete factorization, a multigrid V-cycle
(\cref{ch:multigrid}), or even an inner Krylov solve---because all of them are just
functions of the shape $\vr\mapsto\vz$. Preconditioning is modular for exactly this
reason: choosing a preconditioner is choosing a function, and the iteration is blind to
the choice.
\end{keyidea}

\Cref{fig:blackbox-pc} draws the blindness. GMRES issues an \code{apply} and receives a
vector; whatever algorithm produced that vector is sealed inside the box. Swap the box's
contents---Jacobi for ILU, ILU for multigrid---and not one line of the iteration changes.
The arrow into the box and the vector out of it are the entire contract.

\begin{figure}[t]
\centering
\begin{tikzpicture}[font=\footnotesize, node distance=9mm]
  % the GMRES driver
  \node[stage, minimum width=24mm, minimum height=20mm] (gm) at (0,0)
    {\textbf{GMRES}\\[2pt]\footnotesize the iteration\\asks only\\\code{apply(pc, r)}};
  % the opaque preconditioner box
  \node[draw=simViolet, fill=simBgViolet, very thick, densely dashed, rounded corners=3pt,
        minimum width=30mm, minimum height=30mm, right=20mm of gm] (box) {};
  \node[font=\scriptsize\bfseries, simViolet] at (box.north) [yshift=-3mm] {preconditioner};
  \node[font=\scriptsize\itshape, simGray, align=center] at (box.center) [yshift=2mm]
    {Jacobi?\\ILU?\\multigrid V-cycle?\\inner Krylov solve?\\[2pt]\textbf{GMRES cannot tell.}};
  % apply arrow in
  \draw[flow] (gm.east) -- node[above, font=\scriptsize]{$\vr$}
    node[below, font=\scriptsize]{\code{apply}} (box.west);
  % vector out
  \draw[flow] (box.south) .. controls +(0,-0.9) and +(0,-0.9) ..
    node[below, font=\scriptsize]{$\vz\approx\mat{M}^{-1}\vr$} (gm.south);
  % the seal
  \node[font=\scriptsize\bfseries, simRust, align=center, right=2mm of box]
    {one interface,\\any algorithm};
\end{tikzpicture}
\caption[The preconditioner as a black box GMRES cannot see inside]{Modularity made
visual. The GMRES iteration hands the preconditioner a residual $\vr$ through the single
\code{apply} interface and receives back a vector $\vz\approx\mat{M}^{-1}\vr$. The
algorithm \emph{inside} the box---a diagonal divide, an incomplete factorization, a
multigrid V-cycle, an inner Krylov solve---is completely sealed off: GMRES never inspects
it and never needs to. Replacing the box's contents changes the convergence \emph{rate}
but not a single line of the iteration. This is the solver-as-operator rotation of
\cref{ch:operators}: because an approximate inverse wears the same vector-in, vector-out
shape as an operator, the iteration treats it as one. Choosing a preconditioner is
choosing the function inside the box.}
\label{fig:blackbox-pc}
\end{figure}

\subsection{The same iteration, four preconditioners}

To feel the modularity concretely, here is the preconditioned-residual line of a Krylov
step, written once, with four different bindings of \code{pc}. The iteration code is
identical in all four; only the function handed to it differs.

\begin{lstlisting}[style=l4style]
-- the iteration body is written ONCE against the apply interface:
let z = apply pc r in           -- z ~ M^{-1} r ; pc is opaque here
... use z in the Krylov update ...

-- ...and pc is any of these, chosen at construction, never re-inspected:
let pc = make_jacobi   A         in ...   -- z = d .* r          (one divide)
let pc = make_ilu      A         in ...   -- z = U \ (L \ r)     (two triangular solves)
let pc = make_multigrid hierarchy in ...  -- z = v_cycle r       (a recursive sweep)
let pc = \r -> gmres_inner A r 3  in ...   -- z = 3 steps of inner GMRES (flexible!)
\end{lstlisting}

The first three are constructed operators (\cref{sec:constructed}): build once, apply many
times, the \emph{same} \code{pc} every step. The fourth---an inner solve---produces a
\emph{different} effective operator each step, and that difference, as we will see in
\cref{sec:flexible}, is precisely what forces FGMRES.

\section{The constructed-operator view: build once, apply each step}
\label{sec:constructed}

The preconditioner apply is cheap per call, but most preconditioners are \emph{not} cheap
to set up: extracting and inverting a diagonal, computing an incomplete factorization, or
assembling a multigrid hierarchy is real work. The saving grace is that this work is done
\emph{once}, before the iteration begins, and then amortized over every step. This is
exactly the build-time / run-time phase split of \cref{ch:strata} and the constructed
operator of \cref{ch:operators}.

\begin{keyidea}
A preconditioner is a \emph{constructed operator} (\cref{ch:operators}). Its life has two
phases in two strata. \textbf{Construction} (build-time, once): consume $\mat{A}$ and a
configuration, do the expensive setup---extract a diagonal, compute an incomplete
factorization, build a multigrid hierarchy---and freeze the result into immutable internal
state. \textbf{Application} (run-time, many times): consume a residual $\vr$, read the
frozen state, return $\vz\approx\mat{M}^{-1}\vr$, change nothing. The expensive part runs
once; the cheap part runs every step. This is why the per-iteration surcharge of
\cref{sec:trade} is just \emph{one apply}, not a fresh setup.
\end{keyidea}

\Cref{fig:build-apply} draws the split for a generic preconditioner, deliberately echoing
the construct-then-apply figure of \cref{ch:operators}. The matrix and config flow in on
the left and are consumed once by the build; the opaque preconditioner value sits in the
middle holding its frozen internals; a stream of cheap \code{apply(r)} calls flows out on
the right, one per Krylov step, each reading the same internals.

\begin{figure}[t]
\centering
\begin{tikzpicture}[node distance=10mm]
  % construction inputs
  \node[data, align=left, text width=24mm] (cfg)
    {operator $\mat{A}$\\pc-type enum\\config};
  % build step
  \node[stage, minimum width=26mm, minimum height=14mm, right=12mm of cfg] (build)
    {\textbf{build}\\[1pt]\footnotesize extract diag /\\factorize /\\set up hierarchy};
  \draw[flow] (cfg) -- node[above,font=\scriptsize]{once} (build);
  % opaque pc value
  \node[draw=simViolet, fill=simBgViolet, very thick, densely dashed, rounded corners=3pt,
        minimum width=24mm, minimum height=20mm, right=12mm of build] (pc) {};
  \node[font=\ttfamily\scriptsize, simInk] at (pc.north) [yshift=-3.2mm] {pc};
  \node[font=\scriptsize\itshape, simGray, align=center] at (pc.center)
    {frozen $\mat{M}^{-1}$\\internals\\(diag, factors,\\hierarchy)};
  \draw[flow] (build) -- (pc);
  % many cheap applies
  \foreach \y/\lbl in {1.1/a, 0.0/b, -1.1/c}{
    \node[draw=simTeal, fill=simBgGreen, semithick, rounded corners=2pt,
          minimum width=20mm, minimum height=6mm] (ap\lbl) at ($(pc.east)+(2.7,\y)$)
          {\footnotesize\ttfamily apply r};
    \draw[flow] (pc.east) -- (ap\lbl.west);
  }
  \node[font=\scriptsize\itshape, simGray] at ($(pc.east)+(2.7,-1.9)$)
    {one per Krylov step};
  % stratum band labels
  \node[font=\scriptsize\bfseries, simRust] at ($(build)+(0,-1.7)$) {build-time};
  \node[font=\scriptsize\bfseries, simGreen] at ($(pc.east)+(2.7,1.9)$) {run-time};
\end{tikzpicture}
\caption[Build once, apply each step]{A preconditioner is the constructed operator of
\cref{ch:operators}, with the two phases in the two strata of \cref{ch:strata}.
\emph{Construction} (left) consumes the operator $\mat{A}$, a preconditioner-type enum,
and configuration, runs the expensive setup \emph{once}---extract and invert a diagonal,
compute an incomplete factorization, assemble a multigrid hierarchy---and freezes the
result into the preconditioner's immutable internals. \emph{Application} (right) consumes
only a residual and runs \emph{many} times, once per Krylov step, each call cheap and pure
on the frozen internals. The build-time arrow fires once; the run-time arrows fire
hundreds of times. This phase split is why the per-step cost of preconditioning is a
single apply, never a fresh setup.}
\label{fig:build-apply}
\end{figure}

\subsection{The factory that returns a uniform apply}

A real solver must choose \emph{which} preconditioner to build, and that choice arrives as
a configuration enum. \Cref{ch:operators} named the place the enum is consumed: the
\emph{constructed-operator factory}. A single function reads the
preconditioner-type$\,\times\,$side$\,\times\,$multigrid axes, dispatches once, and returns
a \emph{uniformly typed} apply---past which no enum survives. Palace's
\code{ConfigurePreconditionerSolver} is exactly this factory: it reads the pc-type enum
(\code{JACOBI}, \code{AMS}, \code{BoomerAMG}, a direct factorization, \dots), wraps the
result in a multigrid solver when the FE-space hierarchy has more than one level, and hands
back one \code{Solver} value the iteration applies without ever re-reading the enum.

\begin{lstlisting}[style=l4style]
-- the factory: ONE site reads the pc-type enum, returns a uniform apply
configure_pc :: PcType -> Config -> Op[Tensor[N] -> Tensor[N]]
-- downstream the enum is GONE; only the uniform constructed operator remains:
let pc = configure_pc pc_type cfg in   -- dispatch happens HERE, once
... each Krylov step ...
let z = apply pc r in                  -- ...and never again; pc is opaque
\end{lstlisting}

\begin{workedexample}{Jacobi-preconditioned CG versus plain CG}{jacobi-cg}
We make the whole chapter concrete on the cleanest possible preconditioner---Jacobi
(diagonal) scaling---and a model ill-conditioned operator, and watch the iteration count
collapse.

\textbf{The model problem.} Take a diagonally-dominant SPD operator whose diagonal entries
span a wide range, so that the diagonal carries most of the ill-conditioning. A concrete
small instance: let $\mat{A}$ have diagonal $\operatorname{diag}(\mat{A})=[1,10,100,1000]$
with modest off-diagonal coupling, so $\kappa(\mat{A})\approx 10^{3}$. The spread is
\emph{in the diagonal}---exactly the case Jacobi is built for.

\textbf{Construction (once).} The Jacobi preconditioner is $\mat{M}=\operatorname{diag}(\mat{A})$,
and its apply is ``divide each component by the diagonal entry.'' Construction extracts the
diagonal and inverts it componentwise, producing one vector $\mathbf{d}$ of reciprocals
(\cref{wex:jacobi} of \cref{ch:operators} traced this end to end):
\begin{lstlisting}[style=l4style]
make_jacobi A =
  let d = reciprocal (extract_diagonal A) in   -- O(N), once
  \r -> d .* r                                  -- the apply: O(N) per call
\end{lstlisting}
For our $\mat{A}$, $\mathbf{d}=[1,\tfrac1{10},\tfrac1{100},\tfrac1{1000}]$.

\textbf{The spectral effect.} Scaling each row by its diagonal pulls the four widely-spread
diagonal entries to the \emph{same} value $1$, leaving only the off-diagonal coupling to
perturb the spectrum. The eigenvalues of $\mat{M}^{-1}\mat{A}=\operatorname{diag}(\mat{A})^{-1}\mat{A}$
cluster tightly near $1$: the operator now has unit diagonal, and for our modest coupling
$\kappa(\mat{M}^{-1}\mat{A})\approx 2$ to $3$ instead of $10^3$. This is exactly the
top-to-bottom collapse of \cref{fig:spectrum}.

\textbf{The iteration count.} Feed both condition numbers into the CG bound
$2\big((\sqrt\kappa-1)/(\sqrt\kappa+1)\big)^k$ of \cref{thm:converge}. To reach a relative
energy-norm error of $10^{-6}$:
\begin{itemize}[nosep]
  \item plain CG, $\kappa=10^{3}$: the rate factor is $(\sqrt{1000}-1)/(\sqrt{1000}+1)\approx
  0.939$, so $0.939^k\le 10^{-6}/2$ needs $k\approx 240$ iterations;
  \item Jacobi-CG, $\kappa=2.5$: the rate factor is $(\sqrt{2.5}-1)/(\sqrt{2.5}+1)\approx
  0.226$, so $0.226^k\le 10^{-6}/2$ needs $k\approx 11$ iterations.
\end{itemize}
A factor of roughly $20$ fewer iterations, bought by one $O(N)$ construction and one extra
$O(N)$ multiply per step---a per-step surcharge dwarfed by the operator apply on any real
sparse $\mat{A}$. \Cref{fig:jacobi-conv} plots the two convergence curves; the gap is the
entire point of preconditioning.
\end{workedexample}

\begin{figure}[t]
\centering
\begin{tikzpicture}
  \begin{semilogyaxis}[
      width=12.4cm, height=6.6cm,
      xlabel={iteration $k$}, ylabel={relative energy-norm error (bound)},
      xmin=0, xmax=260, ymin=1e-8, ymax=3,
      grid=both, grid style={simGray!18},
      legend pos=north east, legend cell align=left,
      legend style={font=\footnotesize, draw=simGray!40},
      label style={font=\footnotesize}, tick label style={font=\scriptsize},
      every axis plot/.append style={very thick}]
    % plain CG: kappa = 1000
    \addplot[simRust, domain=0:260, samples=80]
      {2*exp(x*ln((sqrt(1000)-1)/(sqrt(1000)+1)))};
    \addlegendentry{plain CG, $\kappa=10^3$ ($\approx240$ its)}
    % Jacobi-CG: kappa = 2.5
    \addplot[simGreen, domain=0:30, samples=40]
      {2*exp(x*ln((sqrt(2.5)-1)/(sqrt(2.5)+1)))};
    \addlegendentry{Jacobi-CG, $\kappa=2.5$ ($\approx11$ its)}
    \draw[simGray, densely dashed] (axis cs:0,1e-6)--(axis cs:260,1e-6)
      node[pos=0.5, above, font=\scriptsize, text=simGray] {tol $=10^{-6}$};
  \end{semilogyaxis}
\end{tikzpicture}
\caption[Jacobi-preconditioned CG versus plain CG]{The convergence bound of
\cref{thm:converge} for the model problem of \cref{wex:jacobi-cg}. \emph{Plain CG} (rust)
on the bare operator $\kappa=10^3$ crawls, needing about $240$ iterations to reach a
$10^{-6}$ tolerance. \emph{Jacobi-preconditioned CG} (green), running on
$\mat{M}^{-1}\mat{A}$ with its spectrum pulled to $\kappa\approx2.5$, reaches the same
tolerance in about $11$ iterations. The two curves run the \emph{same} CG iteration; the
only difference is that the green one applies the diagonal preconditioner once per step.
The $\approx20\times$ reduction in iterations is bought by one $O(N)$ construction and one
$O(N)$ apply per step---a trade that pays handsomely whenever the diagonal carries the
ill-conditioning.}
\label{fig:jacobi-conv}
\end{figure}

\subsection{When the preconditioner varies: the decision rule and FGMRES}
\label{sec:flexible}

The constructed-operator picture has one precondition: there must be a \emph{single}
$\mat{M}$ to freeze at construction. \Cref{ch:operators} gave the rule that decides when
this holds.

\begin{keyidea}
\textbf{The decision rule (\cref{ch:operators}), applied to preconditioners.} A
preconditioner \emph{fixed} for the whole solve---chosen at solve start, the same operator
every step---is build-time configuration: freeze it into a constructed operator, apply the
same \code{pc} uniformly. A preconditioner that \emph{varies per step}---a different
effective operator $\mat{M}_k$ on each iteration---cannot be frozen, because at construction
time the later $\mat{M}_k$ do not yet exist. It must instead be threaded through the run-time
state. This single rule is exactly why a flexible preconditioner forces FGMRES rather than
GMRES (\cref{ch:gmres}).
\end{keyidea}

When does a preconditioner vary per step? Whenever the apply is \emph{itself an inner
iterative solve} run to a loose, varying tolerance: a few steps of an inner GMRES, a
multigrid cycle whose accuracy drifts, a smoother that adapts to the residual. Each inner
solve produces a slightly different linear operator, so $\vz_k=\mat{M}_k^{-1}\vr_k$ uses a
different $\mat{M}_k$ every step. There is no single $\mat{M}$ to construct---the fourth
binding in the listing of \cref{sec:operator}, \lstinline[style=l4style]|\r -> gmres_inner A r 3|,
is exactly this case.

\Cref{wex:fixed-vs-flexible} contrasts the two apply signatures and shows the variant
moving out of the constructed operator and into the threaded state, exactly as the decision
rule predicts.

\begin{workedexample}{The decision rule in action: a fixed apply versus a per-step apply}{fixed-vs-flexible}
We put two preconditioner applies side by side and read their signatures.

\textbf{Case A: a fixed preconditioner (constructed once).} The preconditioner is built
before the loop and applied---the \emph{same} operator---every step. Its apply signature
takes only the residual; everything else is frozen inside:
\begin{lstlisting}[style=l4style]
let pc = configure_pc JACOBI cfg in    -- built ONCE; pc :: Op[Tensor[N] -> Tensor[N]]
... each step ...
let z = apply pc r in ...               -- apply :: Op[N->N] -> Tensor[N] -> Tensor[N]
                                        -- SAME pc every step; nothing threaded
\end{lstlisting}
The state the loop threads is the Krylov basis and bookkeeping; the preconditioner is
\emph{not} part of it, because it never changes. The constructed operator absorbs the entire
``which preconditioner'' question. This is the left branch of the decision rule
(\cref{fig:decisionrule} of \cref{ch:operators}), and plain GMRES or CG handles it.

\textbf{Case B: a per-step preconditioner (an inner solve).} Now each step runs a fresh
inner solve of varying accuracy. There is no \code{pc} to construct; the apply takes the
\emph{step index} (or the residual that determines the inner tolerance) and produces a
different effective operator each time:
\begin{lstlisting}[style=l4style]
-- NO single pc; M_j is produced per step by an inner solve:
let z_j = apply_M_step j v_j in   -- z_j = M_j^{-1} v_j, a DIFFERENT M_j each j
\end{lstlisting}
and---the crucial consequence---the per-step preconditioned vectors $\vz_j$ can no longer be
recomputed on demand, because the inner solve that made each one is gone by the next step.
They must be \emph{remembered}. The solver threads a \emph{second} basis $\mat{Z}_m=[\vz_1\
\cdots\ \vz_m]$ alongside the Arnoldi basis $\mat{V}_m$, and recovers its iterate from the
remembered preconditioned basis:
\begin{lstlisting}[style=l4style]
-- GMRES  (fixed pc):    x = x_0 + V_m y_m   -- pc is a constructed operator, NOT in state
-- FGMRES (per-step pc): x = x_0 + Z_m y_m   -- the per-step Z basis IS threaded state
\end{lstlisting}

\textbf{The reading.} Nothing went wrong; a real distinction surfaced. The fixed
preconditioner of Case A is build-time configuration---freeze it, apply uniformly. The
per-step preconditioner of Case B is run-time state---thread it, and the carry grows by one
extra basis $\mat{Z}$. The decision rule read the variant's lifetime and routed it to the
matching stratum. ``FGMRES exists because the preconditioner became per-step state'' is the
whole story, told in full in \cref{ch:gmres}.
\end{workedexample}

\begin{pitfall}
\textbf{Do not feed a per-step preconditioner to plain GMRES.} Plain GMRES recovers its
iterate as $\vx_0+\mat{M}^{-1}\mat{V}_m\vy$, applying one $\mat{M}^{-1}$ to the whole
correction under the assumption of a \emph{single fixed} $\mat{M}$. Hand it a varying
$\mat{M}_k$ and the code still runs---no crash---and even reports a small residual from its
running QR; but that residual belongs to a least-squares problem the reconstructed iterate
no longer matches, because the reconstruction used the wrong inverse. The solve converges to
the \emph{wrong answer}, silently. The fix is not a patch to GMRES; it is FGMRES, which
remembers the preconditioned vectors in the $\mat{Z}$ basis (\cref{ch:gmres}). The tell that
you need it: your preconditioner apply is itself an iterative solve.
\end{pitfall}

\section{A tour of preconditioner families}
\label{sec:families}

Every preconditioner is a function $\vr\mapsto\vz\approx\mat{M}^{-1}\vr$; the families
differ only in \emph{which} function, i.e. in what their construction freezes and what their
apply computes. We survey the main families here as apply-functions, light on internals---each
gets its own chapter later. The point is the unity: they all wear the same interface, so
choosing one is choosing a function, and the iteration is unchanged.

\begin{description}[style=nextline,leftmargin=1.5em]
  \item[Diagonal / Jacobi.] $\mat{M}=\operatorname{diag}(\mat{A})$; the apply is a single
  componentwise divide $\vz=\mathbf{d}\odot\vr$. The cheapest possible preconditioner, and
  the constructed-operator archetype (\cref{wex:jacobi-cg}). It clusters well only when the
  diagonal carries the ill-conditioning; it is also the simplest \emph{smoother}
  (\cref{ch:smoothers}).
  \item[Incomplete factorizations (ILU / IC).] Compute an approximate $\mat{L}\mat{U}$ (or,
  for SPD $\mat{A}$, Cholesky $\mat{L}\mat{L}^{\top}$) factorization that keeps only a
  sparsity pattern, discarding fill-in. Construction is the incomplete factorization; the
  apply is two triangular solves $\vz=\mat{U}^{-1}(\mat{L}^{-1}\vr)$. Stronger than Jacobi,
  more expensive to build and apply, and---being inherently sequential---a poor fit for the
  GPU target (\cref{ch:strata} discusses the recurrence obstruction).
  \item[Domain decomposition.] Partition the mesh into subdomains, solve a smaller problem on
  each, and stitch the pieces together. The apply is a collection of independent subdomain
  solves plus a coarse correction. Naturally parallel across subdomains; the family of choice
  when the problem is already distributed.
  \item[Multigrid.] The big one for elliptic PDEs (\cref{ch:multigrid}). Construction builds a
  \emph{hierarchy} of coarser and coarser meshes; the apply is a V-cycle that smooths the
  high-frequency error on fine levels (with a cheap smoother like Jacobi) and solves the
  low-frequency error on coarse levels (cheaply, because they are small). Done right, multigrid
  is \emph{optimal}: its cluster, and hence its iteration count, is independent of mesh size
  (\cref{sec:meshindep}).
  \item[Auxiliary-space / AMS.] For $\Hcurl$ curl--curl problems (the magnetostatic and
  eigenmode operators), the natural multigrid smoothers fail on the large near-kernel of the
  curl operator. The \emph{auxiliary-space Maxwell solver} (AMS) preconditions by mapping the
  problem into auxiliary $\Hone$ scalar and vector spaces---where ordinary multigrid works---and
  mapping back. Its apply is a fixed composition of grid-transfer and smoothing operations
  (\cref{ch:smoothers}).
\end{description}

\Cref{fig:family-tree} arranges the families as a tree, all sharing one root: the
\code{apply} interface. The branches are different functions; the root is the single shape
$\vr\mapsto\vz$ that lets the same Krylov method drive any of them.

\begin{figure}[t]
\centering
\begin{tikzpicture}[font=\footnotesize,
    famnode/.style={draw=simViolet, fill=simBgViolet, rounded corners=2pt, semithick,
      align=center, inner sep=3pt, minimum height=8mm, text width=22mm},
    node distance=4mm]
  % root: the apply interface
  \node[draw=simRust, fill=simBgRust, very thick, rounded corners=3pt, align=center,
        inner sep=4pt, minimum width=46mm]
    (root) at (0,3.2)
    {\textbf{the \code{apply} interface}\\[1pt]\footnotesize $\vr\mapsto\vz\approx\mat{M}^{-1}\vr$};
  % the five families
  \node[famnode] (jac) at (-5.4,0.8) {Jacobi /\\diagonal\\\scriptsize one divide};
  \node[famnode] (ilu) at (-2.7,0.8) {ILU / IC\\\scriptsize triangular\\\scriptsize solves};
  \node[famnode] (dd)  at (0,0.8)    {domain\\decomposition\\\scriptsize subdomain solves};
  \node[famnode] (mg)  at (2.7,0.8)  {multigrid\\\scriptsize V-cycle\\\scriptsize (optimal)};
  \node[famnode] (ams) at (5.4,0.8)  {AMS\\\scriptsize aux-space\\\scriptsize ($\Hcurl$)};
  % connect root to each
  \foreach \n in {jac,ilu,dd,mg,ams}{
    \draw[depedge] (root.south) -- (\n.north);}
  % chapter pointers
  \node[font=\scriptsize\itshape, simGray, below=1mm of jac] {\cref{ch:smoothers}};
  \node[font=\scriptsize\itshape, simGray, below=1mm of mg]  {\cref{ch:multigrid}};
  \node[font=\scriptsize\itshape, simGray, below=1mm of ams] {\cref{ch:smoothers}};
\end{tikzpicture}
\caption[A family tree of preconditioners sharing one interface]{The preconditioner
families, arranged as a tree rooted in the single \code{apply} interface
$\vr\mapsto\vz\approx\mat{M}^{-1}\vr$. Jacobi/diagonal scaling, incomplete factorizations
(ILU/IC), domain decomposition, multigrid, and the auxiliary-space AMS solver differ
entirely in \emph{what} their apply computes and \emph{what} their construction freezes---a
divide, a pair of triangular solves, a set of subdomain solves, a V-cycle, a mapping into
auxiliary spaces---but every one of them is a function of the same shape. Because they share
the interface, the same Krylov method drives any of them, and choosing a preconditioner is
choosing which branch of the tree to hang in the black box of \cref{fig:blackbox-pc}. The
chapters named below the leaves develop the internals.}
\label{fig:family-tree}
\end{figure}

\section{Spectral equivalence and mesh-independent convergence}
\label{sec:meshindep}

We close with the property that separates a merely good preconditioner from an
\emph{optimal} one, and that is the real goal for the elliptic operators Palace solves.

The condition-number improvement of \cref{wex:jacobi-cg} was a fixed factor: $\kappa$ from
$10^3$ to $2.5$ on \emph{that} mesh. But Palace refines meshes, and the question that
matters is what happens as the mesh shrinks. For the bare elliptic operators, the condition
number \emph{grows} as the mesh is refined: discretizing $-\Div(\varepsilon\Grad\,\cdot\,)$
on a mesh of element size $h$ gives $\kappa(\mat{A})=O(h^{-2})$, so halving $h$ quadruples
$\kappa$ and---since CG iterations scale like $\sqrt\kappa$---roughly \emph{doubles} the
iteration count. A solver whose iteration count climbs every time you refine is a solver
that does not scale.

\begin{definition}[Spectral equivalence and optimality]
\label{def:specequiv}
A preconditioner $\mat{M}$ is \emph{spectrally equivalent} to $\mat{A}$ if there are
constants $0<c_1\le c_2$, \emph{independent of the mesh size $h$}, such that
\[
  c_1\,\vx^{\top}\mat{M}\vx \;\le\; \vx^{\top}\mat{A}\vx \;\le\; c_2\,\vx^{\top}\mat{M}\vx
  \qquad\text{for all }\vx.
\]
Then the preconditioned operator satisfies $\kappa(\mat{M}^{-1}\mat{A})\le c_2/c_1$,
\emph{bounded independently of $h$}. A preconditioner achieving this is called
\emph{optimal}: the iteration count to a fixed tolerance does not grow as the mesh is
refined.
\end{definition}

\begin{keyidea}
The real goal of preconditioning an elliptic PDE is not a one-off condition-number
improvement---it is \emph{mesh independence}. A spectrally-equivalent (optimal)
preconditioner makes $\mat{M}^{-1}\mat{A}$ as well-conditioned on a fine mesh as on a coarse
one: the cluster of \cref{fig:spectrum} stays the same width no matter how many unknowns
$N$. The iteration count then stays \emph{flat} as $N$ grows, so the total solve cost scales
like $O(N)$---linearly---rather than $O(N^{1.5})$ or worse. Multigrid (\cref{ch:multigrid})
is the canonical optimal preconditioner for these operators, and mesh-independent
convergence is exactly the property that makes it indispensable at scale.
\end{keyidea}

\Cref{fig:meshindep} contrasts the two regimes: a fixed-strength preconditioner (Jacobi),
whose iteration count climbs as $h\to0$ because $\kappa$ still grows like $h^{-2}$ underneath
it, against an optimal preconditioner (multigrid), whose iteration count is a flat line
because the cluster width is bounded. The flat line is the prize.

\begin{figure}[t]
\centering
\begin{tikzpicture}
  \begin{axis}[
      width=12.4cm, height=6.6cm,
      xlabel={mesh refinement $1/h$ (finer mesh $\to$)},
      ylabel={iterations to tolerance},
      xmin=0, xmax=34, ymin=0, ymax=130,
      tick label style={font=\scriptsize}, label style={font=\footnotesize},
      legend pos=north west, legend cell align=left,
      legend style={font=\footnotesize, draw=simGray!40},
      grid=both, grid style={simGray!15},
      every axis plot/.append style={very thick, mark=*, mark size=1.4pt}]
    % Jacobi / fixed preconditioner: iterations ~ sqrt(kappa) ~ sqrt(h^-2) = 1/h, linear climb
    \addplot[simRust] coordinates
      {(2,8)(4,16)(8,32)(12,48)(16,64)(20,80)(24,96)(28,112)(32,124)};
    \addlegendentry{Jacobi (mesh-dependent): iters $\sim 1/h$}
    % multigrid / optimal: flat
    \addplot[simGreen] coordinates
      {(2,9)(4,9)(8,10)(12,10)(16,10)(20,11)(24,11)(28,11)(32,11)};
    \addlegendentry{multigrid (optimal): iters bounded}
  \end{axis}
\end{tikzpicture}
\caption[Mesh-dependent versus mesh-independent convergence]{Iteration count as the mesh is
refined ($1/h$ growing rightward). With a \emph{fixed-strength} preconditioner like Jacobi
(rust), the underlying condition number still grows like $h^{-2}$, so the iteration count
climbs roughly like $1/h$---refine the mesh and every solve gets slower, the hallmark of a
non-scalable method. With an \emph{optimal}, spectrally-equivalent preconditioner like
multigrid (green, \cref{def:specequiv}), the condition number of $\mat{M}^{-1}\mat{A}$ is
bounded independently of $h$, so the iteration count is essentially \emph{flat} no matter how
fine the mesh. The flat green line is the goal of preconditioning at scale: total cost
$O(N)$, not $O(N^{1.5})$. Mesh independence is what makes multigrid (\cref{ch:multigrid})
indispensable for the elliptic operators Palace assembles.}
\label{fig:meshindep}
\end{figure}

\subsection{Tie-off: preconditioned CG and the SPD requirement}
\label{sec:pcg-tie}

\Cref{ch:cg} already showed (\cref{sec:pcg}) how the CG \emph{loop} changes under
preconditioning: form the preconditioned residual $\vz=\mat{M}^{-1}\vr$ once per step, and
replace the residual inner products $\inner{\vr}{\vr}$ in $\alpha$ and $\beta$ by the
$\mat{M}$-weighted products $\inner{\vz}{\vr}$. We can now name \emph{why} the loop takes that
particular shape. Running CG naively on $\mat{M}^{-1}\mat{A}$ would break SPD-ness, because
$\mat{M}^{-1}\mat{A}$ is generally not symmetric (\cref{sec:lrs}); preconditioned CG instead
runs CG in the inner product defined by $\mat{M}$, which is precisely the split-preconditioning
view $\mat{L}^{-1}\mat{A}\mat{L}^{-\top}$ for $\mat{M}=\mat{L}\mat{L}^{\top}$. That operator
\emph{is} SPD, so every guarantee of \cref{ch:cg} survives. The discipline this forces---the
preconditioner $\mat{M}$ must itself be SPD for preconditioned CG to stay honest---is the
pitfall \cref{ch:cg} flagged, now explained: it is the price of keeping the energy inner
product an inner product.

\summarysection
A \emph{preconditioner} $\mat{M}\approx\mat{A}$, cheap to invert, transforms
$\mat{A}\vx=\vb$ into an equivalent system whose operator $\mat{M}^{-1}\mat{A}$ has a
spectrum clustered near $1$, so the \emph{same} Krylov iteration (\cref{ch:cg,ch:gmres})
converges in far fewer steps. The fundamental trade is a product: each step costs one extra
preconditioner apply, but the iteration count falls, and the win comes when $\mat{M}^{-1}$ is
\emph{both} a good approximation of $\mat{A}^{-1}$ (clusters the spectrum) and cheap to apply.
The preconditioner can be placed on the \emph{left} ($\mat{M}^{-1}\mat{A}\vx=\mat{M}^{-1}\vb$,
distorts the residual), the \emph{right} ($\mat{A}\mat{M}^{-1}\vu=\vb$, preserves the true
residual---the GMRES default), or \emph{split} ($\mat{M}_1^{-1}\mat{A}\mat{M}_2^{-1}$,
preserves SPD-ness for CG)---a variant axis the construction absorbs. The chapter's spine is
that \emph{a preconditioner is an operator you apply}: $\mat{M}^{-1}$ is never formed, only
applied as a black-box function $\vr\mapsto\vz\approx\mat{M}^{-1}\vr$, and the Krylov method
needs nothing else---so the same GMRES drives Jacobi, ILU, multigrid, or an inner solve
without changing a line (\cref{ch:operators}'s solver-as-operator move). The preconditioner is
a \emph{constructed operator}: build once (extract a diagonal, factorize, set up a hierarchy)
at build-time, apply cheaply each step at run-time. When the preconditioner is \emph{fixed},
the constructed-operator factory absorbs it; when it \emph{varies per step} (an inner solve),
the decision rule routes it into the threaded state and forces FGMRES. The families---Jacobi,
ILU/IC, domain decomposition, multigrid, auxiliary-space AMS---all share the apply interface;
choosing one is choosing a function. The real goal for elliptic operators is \emph{spectral
equivalence}: an optimal preconditioner bounds $\kappa(\mat{M}^{-1}\mat{A})$ independently of
the mesh size $h$, so the iteration count stays flat as the mesh is refined and the solve
scales linearly---the property that makes multigrid (\cref{ch:multigrid}) indispensable.

\furtherreading
The definitive treatment of preconditioned Krylov methods---left/right/split placement, the
flexible-preconditioner distinction that drives FGMRES, and the spectral-equivalence theory of
optimal preconditioners---is \citet{saad2003} (Chapters~9 and~10), which also motivates the
black-box ``a Krylov method needs only the action $\vv\mapsto\mat{M}^{-1}\vv$'' abstraction
that makes preconditioning modular. \citet{trefethenbau1997} gives the cleanest short account
of how clustering the spectrum---rather than merely shrinking $\kappa$---governs Krylov
convergence (Lectures~35 and~40). For the matrix-computation mechanics of incomplete
factorizations and the diagonal/Jacobi family, \citet{golubvanloan2013} is the standard
reference. The optimal multigrid preconditioner and the auxiliary-space AMS solver for
$\Hcurl$ are developed in \cref{ch:multigrid} and \cref{ch:smoothers}; the constructed-operator
and solver-as-operator abstractions the chapter rests on are \cref{ch:operators}.

\exercisessection
\begin{enumerate}[nosep]
  \item \textbf{The two demands.} \Cref{def:precond} asks a preconditioner to be both a good
  approximation of $\mat{A}^{-1}$ and cheap to apply. Explain why the two extremes
  $\mat{M}=\mat{I}$ and $\mat{M}=\mat{A}$ each satisfy exactly one demand and fail the other,
  and what each does to (i) the iteration count and (ii) the per-step cost. Where on the trade
  curve of \cref{fig:cost-trade} does each extreme sit?
  \item \textbf{The trade, quantified.} A solve takes $200$ iterations unpreconditioned, each
  costing one operator apply of $C$ units. A preconditioner cuts the iteration count to $25$ but
  adds a preconditioner apply costing $0.8C$ per step. (i) Compute the total cost with and
  without preconditioning (ignore the one-time build). (ii) Now suppose a \emph{stronger}
  preconditioner cuts iterations to $8$ but its apply costs $5C$ per step; compute its total
  cost. (iii) Which of the three wins, and what does this illustrate about \cref{fig:cost-trade}?
  \item \textbf{Clustering versus $\kappa$.} \Cref{ch:cg} noted that CG's true convergence depends
  on the spectrum's \emph{layout}, not just $\kappa$. Two preconditioned operators both have
  $\kappa=50$: operator $P$ has its eigenvalues spread uniformly across $[1,50]$, operator $Q$ has
  $48$ eigenvalues tightly clustered near $1$ and two outliers near $50$. Which converges faster
  under CG, and why? Relate your answer to the polynomial picture of \cref{sec:cluster} and to
  \cref{fig:spectrum}.
  \item \textbf{Left versus right.} (i) Write the system, the operator the iteration runs on, and
  the residual the iteration sees, for left and for right preconditioning. (ii) Explain precisely
  why right preconditioning is preferred for GMRES, referring to the free residual readout of
  \cref{ch:gmres}. (iii) For an SPD $\mat{A}$ with $\mat{M}=\mat{L}\mat{L}^{\top}$, show that the
  split-preconditioned operator $\mat{L}^{-1}\mat{A}\mat{L}^{-\top}$ is symmetric, and explain why
  CG needs this where GMRES does not.
  \item \textbf{The black box.} The same GMRES drives a Jacobi preconditioner, an ILU, and a
  multigrid V-cycle without changing a line (\cref{fig:blackbox-pc}). (i) State the single
  interface that makes this possible and exactly what GMRES asks of the preconditioner through it.
  (ii) Name one thing GMRES is \emph{prevented} from seeing about the preconditioner, and explain
  why that prevention is a feature. (iii) Which chapter's abstraction (\cref{ch:operators}) is this
  an instance of, and what is the shared shape that lets a solver wear an operator's interface?
  \item \textbf{Build once, apply many.} For the Jacobi preconditioner of \cref{wex:jacobi-cg},
  state the cost of construction and of a single apply, each in terms of $N$. If a solve runs $k$
  preconditioned steps, write the total preconditioner work as a function of $N$ and $k$ and
  identify the amortized construction term. For what range of $k$ is constructing the operator
  clearly worthwhile?
  \item \textbf{Apply the decision rule.} For each preconditioner below, decide whether it is a
  constructed operator (fixed, build-time) or per-step threaded state (forcing FGMRES), and
  justify with the decision rule of \cref{sec:flexible}: (a) a diagonal Jacobi scaling built from
  $\mat{A}$ at solve start; (b) an incomplete Cholesky factorization computed once; (c) ``three
  steps of an inner GMRES'' run afresh each outer step; (d) a multigrid V-cycle with fixed smoother
  counts; (e) a multigrid cycle whose inner tolerance adapts to the current residual.
  \item \textbf{Mesh independence.} The bare elliptic operator has $\kappa(\mat{A})=O(h^{-2})$.
  (i) Using the CG iteration estimate iterations $\sim\sqrt{\kappa}$, show that halving $h$ roughly
  doubles the unpreconditioned iteration count. (ii) Define spectral equivalence
  (\cref{def:specequiv}) and explain why it makes the preconditioned iteration count independent of
  $h$. (iii) Sketch, from \cref{fig:meshindep}, the iteration-count-versus-$1/h$ curve for a
  mesh-dependent and an optimal preconditioner, and state the total-cost scaling ($O(N)$ versus
  worse) each implies.
  \item \textbf{The SPD trap.} Preconditioned CG requires the preconditioner $\mat{M}$ to be SPD.
  (i) Explain, using \cref{sec:pcg-tie}, why running plain CG on $\mat{M}^{-1}\mat{A}$ with a
  general $\mat{M}$ breaks the method even when $\mat{A}$ is SPD. (ii) Show how the split view
  $\mat{L}^{-1}\mat{A}\mat{L}^{-\top}$ restores symmetry. (iii) What goes wrong, observably, if a
  non-SPD preconditioner is used inside preconditioned CG?
\end{enumerate}
